\documentclass{article}
\usepackage{enumitem}
\usepackage{amsmath}
\begin{document}

\title{Chapter 25: Human Reproduction}
\date{}
\maketitle

\begin{enumerate}[label=Q\arabic*.]

\item Spermatogenesis occurs in which structure of the testis?
\\(A) Epididymis
\\(B) Seminiferous tubules
\\(C) Rete testis
\\(D) Vas deferens

\item The acrosome of a sperm contains:
\\(A) Mitochondria for energy production
\\(B) Hydrolytic enzymes (acrosin, hyaluronidase) for penetrating the egg
\\(C) Genetic material (DNA)
\\(D) Flagellum proteins for motility

\item Which hormone stimulates Sertoli cells to produce ABP (androgen-binding protein) in males?
\\(A) LH
\\(B) FSH
\\(C) Testosterone
\\(D) GnRH

\item The corpus luteum, formed after ovulation, produces:
\\(A) FSH and LH
\\(B) Progesterone and estrogen to maintain early pregnancy
\\(C) HCG (human chorionic gonadotropin)
\\(D) Oxytocin

\item Assertion: HCG (human chorionic gonadotropin) is the basis for early pregnancy tests.\\
Reason: HCG is secreted by the trophoblast/placenta soon after implantation and maintains the corpus luteum.
\\(A) Both true, Reason is correct explanation
\\(B) Both true, Reason is not correct explanation
\\(C) Assertion true, Reason false
\\(D) Assertion false, Reason true

\item Implantation of the blastocyst in the uterus occurs approximately how many days after fertilization?
\\(A) 1--2 days
\\(B) 3--4 days
\\(C) 6--7 days
\\(D) 14 days

\item The zona pellucida of the secondary oocyte prevents polyspermy by:
\\(A) Blocking sperm entry physically before fertilization
\\(B) Cortical reaction hardening the zona pellucida after the first sperm enters
\\(C) Releasing sperm-killing enzymes
\\(D) Attracting only one sperm at a time

\item Spermiogenesis refers to:
\\(A) Formation of spermatogonia from primordial germ cells
\\(B) Meiotic division of primary spermatocytes
\\(C) Transformation of spermatids into mature spermatozoa
\\(D) Capacitation of sperm in the female tract

\item The primary oocyte is arrested at which stage of meiosis in the human ovary at birth?
\\(A) Prophase I (diplotene stage)
\\(B) Metaphase I
\\(C) Metaphase II
\\(D) Prophase II

\item The placenta performs all of the following functions EXCEPT:
\\(A) Gas exchange between mother and fetus
\\(B) Nutrient and waste transfer
\\(C) Production of HCG, progesterone, and estrogen
\\(D) Production of red blood cells for the fetus

\item Which hormone surge triggers ovulation in the menstrual cycle?
\\(A) FSH surge
\\(B) LH surge (midcycle, around day 14)
\\(C) Estrogen surge
\\(D) Progesterone surge

\item The three primary germ layers of the gastrula and their derivatives are:
\\(A) Ectoderm (skin, nervous system), Mesoderm (muscle, bone, heart), Endoderm (gut lining, lungs)
\\(B) Ectoderm (gut lining), Mesoderm (skin), Endoderm (muscle)
\\(C) Ectoderm (muscle), Mesoderm (nervous system), Endoderm (skin)
\\(D) All three layers produce all tissue types equally

\item The Fallopian tube (oviduct) is the site of:
\\(A) Implantation
\\(B) Fertilization (in the ampullary region)
\\(C) Meiosis II completion
\\(D) Spermatogenesis

\item Match the following hormones with their roles in female reproduction:\\
1. FSH $\rightarrow$ (a) Stimulates follicle development and estrogen production\\
2. LH $\rightarrow$ (b) Triggers ovulation and corpus luteum formation\\
3. Estrogen $\rightarrow$ (c) Proliferative phase of endometrium\\
4. Progesterone $\rightarrow$ (d) Secretory phase of endometrium, maintains pregnancy
\\(A) 1-a, 2-b, 3-c, 4-d
\\(B) 1-b, 2-a, 3-d, 4-c
\\(C) 1-c, 2-d, 3-a, 4-b
\\(D) 1-d, 2-c, 3-b, 4-a

\item Ectopic pregnancy refers to:
\\(A) Implantation of blastocyst in the uterine wall
\\(B) Implantation occurring outside the uterus, most commonly in the Fallopian tube
\\(C) Development of twins in the uterus
\\(D) Failure of implantation leading to early miscarriage

\end{enumerate}
\end{document}
